% Generated draft: requires mathematical and source review.
\documentclass[11pt]{article}
\usepackage[margin=25mm]{geometry}
\usepackage{fontspec}
\setmainfont{Latin Modern Roman}
\usepackage{amsmath,amssymb,mathtools,unicode-math}
\setmathfont{Latin Modern Math}
\usepackage{enumitem,needspace,xurl,hyperref}
\setlength{\parindent}{0pt}
\setlength{\parskip}{5pt}
\setlength{\emergencystretch}{3em}
\newcommand{\N}{\mathbb{N}}
\newcommand{\Z}{\mathbb{Z}}
\newcommand{\Q}{\mathbb{Q}}
\newcommand{\R}{\mathbb{R}}
\newcommand{\C}{\mathbb{C}}
\newcommand{\F}{\mathbb{F}}
\newcommand{\A}{\mathbb{A}}
\newcommand{\PP}{\mathbb{P}}
\newcommand{\E}{\mathbb{E}}
\newcommand{\eps}{\varepsilon}
\newcommand{\dd}{\,\mathrm{d}}
\begin{document}
\section*{UnsolvedMath ID: 700872}
\textbf{Source problem code:} catalogue entry
\subsection{Definitions and mathematical statement}
An \(\ell\)-group is a group \(G\) equipped with a lattice order \(\leq\) such that \(a\leq b\) implies \(xay\leq xby\) for all \(a,b,x,y\in G\). Write \(a\wedge b\) and \(a\vee b\) for the lattice operations. No commutativity assumption is imposed.

For a set \(X\), the free \(\ell\)-group \(F_{\ell}(X)\) has a distinguished map \(\iota:X\to F_{\ell}(X)\) with the following universal property: every map from \(X\) into an \(\ell\)-group \(H\) extends uniquely to a group homomorphism \(F_{\ell}(X)\to H\) preserving \(\wedge\) and \(\vee\). Elements are represented by finite terms in generators from \(X\), the identity, multiplication, inversion, and the two lattice operations. Two elements \(u,v\in F_{\ell}(X)\) are conjugate if
\[
\exists c\in F_{\ell}(X)\qquad c^{-1}uc=v.
\]

The supplied question asks whether conjugacy in a free \(\ell\)-group is decidable: given two such finite terms representing \(u\) and \(v\), is there an algorithm that always halts and correctly decides whether the displayed condition holds?

An effective specification of the generators is required for this algorithmic formulation. The source does not specify the rank of the free \(\ell\)-group, whether a particular rank is fixed, or whether a single algorithm must handle varying finite ranks. These qualifications remain unresolved; the question is retained without silently selecting one of these scopes.
\subsection{Short English statement}
Given two elements of a free lattice-ordered group, can an algorithm always decide whether one is conjugate to the other? The source leaves the rank and the required uniformity across ranks unspecified. Its designation as open is supplied registry metadata.
\subsection{Sources}
\begin{itemize}
\item Local UnsolvedMath registry: supplied complete source record, permanent ID \(700872\), exact source problem code \texttt{catalogue entry}, upstream revision \texttt{d371ee917ac9647fbb06cd6fc529deec6f3dda82}. The record labels the problem \emph{open}; this label is supplied metadata, not an independently verified mathematical assessment.
\item The supplied background attributes Problem 12.53 to S. McCleary, proposed in 1992, in \emph{The Kourovka Notebook (Unsolved Problems in Group Theory)}, 12th issue (1992). It also identifies the 21st-issue edition (2026), arXiv:1401.0300v48. These bibliographic details are supplied attribution and have not been independently checked.
\item The supplied collection description identifies E. I. Khukhro and V. D. Mazurov as editors of the Kourovka Notebook collection covering Issues 1--21 (1965--2026).
\end{itemize}
\end{document}
